% top_mode_edge.tex -- round 415, the edge balance of the
% canonical hole-top mode (reviewer DCCLXXIX, parallel to R413).
% CHART_IDENTITY_EXACT / TOP_MODE_NOT_DEFECT.
% Builds with pdflatex.
% Companion machine check: verify_top_mode_edge.py.
% Discovery: experiments/tfpt-discovery/top_mode_edge_probe.py.
% NO RH CLAIM.  Research documentation, not a theorem of RH.
\documentclass[11pt]{article}

\usepackage[a4paper,margin=2.7cm]{geometry}
\usepackage{amsmath,amssymb,amsthm}
\usepackage{microtype}
\usepackage{booktabs}
\usepackage{xcolor}
\usepackage[colorlinks=true,linkcolor=blue!50!black,urlcolor=blue!50!black]{hyperref}

\newtheorem{lemma}{Lemma}
\newtheorem{prop}{Proposition}
\theoremstyle{remark}
\newtheorem{remark}{Remark}

\newcommand{\indm}{\operatorname{ind}_{-}}
\newcommand{\Tfrak}{\mathfrak{T}}
\newcommand{\sch}{\mathrm{sch}}

\title{\bfseries The top-mode edge\\[2mm]
\large $-\sch=\beta-\alpha$ is the r401 chart rewrite;
the barycentric hole-top mode is not the bulk defect}
\author{Stefan Hamann}
\date{August 29, 2026}

\begin{document}
\maketitle

\begin{abstract}\noindent
After rounds 407--411 the reviewer (DCCLXXIX) asks that the
edge half of P1 be re-read as a \emph{top-mode balance}:
if the hole-top-mode theorem holds in the bulk, the only
critical direction is the barycentric mode
$v_{\mathrm{top},i}\propto 1/(\sqrt{u^\vee(y_i)}\,P_Y'(y_i))$,
and $-\sch_k=\beta_k-\alpha_k$ with $\alpha$ the defect energy
after bulk continuation and $\beta$ the net CD lift after
border consumption.
Round 405 tested the constant (the wrong vector); the Euler
tail and the disk Parseval identity are kept and re-read on
$v_{\mathrm{top}}$.
The Woodbury--Schur rewrite $-\sch=\beta-\alpha$ is a finite
SATZ, over $\mathbb{Q}$ and at machine precision, once
$\alpha=(\eta-1)/\kappa$ and
$\beta=r^T K_W^{-1}r/\kappa$ with the r405 factor
$\kappa=\Delta/(-\sch)$.
It does not depend on which bulk vector is used as $\ell$:
the living/dead census $\beta>\alpha$ / $\beta<\alpha$ is
exactly $\sch<0$ / $\sch>0$ (r401), for $v_{\mathrm{top}}$
and for the constant.
The formula $v_{\mathrm{top}}$ is \emph{not} the bulk defect
($\cos(v_{\mathrm{top}},v_{\mathrm{SV}})=0.68$ at $k_z=9$,
$0.058$ at $k_z=18$ where $\|\Tfrak_0\|=1.25$).
Euler-on-$v_{\mathrm{top}}$ is the constant content
$|\langle 1,v_{\mathrm{top}}\rangle|^2$, not $\alpha_T$.
No RH claim.
\end{abstract}

\medskip\noindent
\textbf{Status.}\enspace
Lemmas~\ref{lem:tail}--\ref{lem:ab} are proved (finite
identities).
Propositions~\ref{prop:euler}--\ref{prop:census} are finite
machine checks.
Ausgang \texttt{CHART\_IDENTITY\_EXACT / TOP\_MODE\_NOT\_DEFECT}.
No $L^*$ claim.  No $R^\dagger$ claim.  No RH claim.

\section{Definitions}\label{sec:def}

On the hole nodes $Y$ let $u^\vee$ be the dual weight of
r356/r407 and $P_Y(x)=\prod_{y\in Y}(x-y)$.
The barycentric hole-top mode in the $\sqrt{u^\vee}$-weighted
sample coordinate is
\begin{equation}\label{eq:vtop}
  v_{\mathrm{top},i}
  \;\propto\;
  \frac{1}{\sqrt{u^\vee(y_i)}\,P_Y'(y_i)}.
\end{equation}
Write $A_0=R_{N-3}-\tfrac12 I$ on $Y$ and $U$ for the two last
dual CD columns (r367).  The ones-split of r405 with
$\ell=v_{\mathrm{top}}$ (normalised) is
$H=A_0+\ell\ell^T$,
\[
  \eta=\ell^T H^{-1}\ell,\qquad
  r=U^T H^{-1}\ell,\qquad
  K_W=I+U^T H^{-1}U,\qquad
  \Delta=1-\eta+r^T K_W^{-1}r.
\]
Set $\alpha_{\mathrm{raw}}=\eta-1$,
$\beta_{\mathrm{raw}}=r^T K_W^{-1}r$, so
$\Delta=\beta_{\mathrm{raw}}-\alpha_{\mathrm{raw}}$.
The r401 scalar $\sch$ is the last Schur complement of the
mixed $3\times 3$ $\Phi$ (two CD directions plus the
Sherman--Morrison border).
With $\kappa_{\mathrm{emp}}=\Delta/(-\sch)$ the chart-scaled
coordinates are
$\alpha=\alpha_{\mathrm{raw}}/\kappa_{\mathrm{emp}}$ and
$\beta=\beta_{\mathrm{raw}}/\kappa_{\mathrm{emp}}$.
Independently, the T0-Rayleigh
$\alpha_T=v_{\mathrm{top}}^T(\Tfrak_0^*\Tfrak_0-I)v_{\mathrm{top}}$
is the candidate ``defect energy after bulk continuation''.

\section{Finite identities}\label{sec:satz}

\begin{lemma}[Euler tail]\label{lem:tail}
For every $z$ and every integer $K\ge 0$,
$1-z^{K+1}=(1-z)\sum_{r=0}^K z^r$.
Over $\mathbb{Q}$, at $z=1/3$ and $K=5$, both sides equal
$728/729$.
In particular $E_p+T_p=1$ \emph{pointwise} on $Y$ for each
geometric prime factor.
\end{lemma}

\begin{lemma}[disk Parseval]\label{lem:disk}
For $0<q<1$ and $K\ge 0$,
$(1-q)\sum_{r=0}^K q^r=1-q^{K+1}$.
At $q=1/2$ and $K=6$ the reserve $q^{K+1}$ is $1/128$.
\end{lemma}

\begin{lemma}[chart identity]\label{lem:ab}
On the r405 rational $3\times 3$
($H=\mathrm{diag}(2,3,4)$, $\ell=(1,1,1)$, $U=I$,
$J^{-1}=\mathrm{diag}(1,1,2)$) one has
$\eta=13/12$, $\Delta=7/36$, $\sch=21/8$,
$\kappa=-2/27$, and with the definitions of
\S\ref{sec:def}
\[
  \alpha=-\tfrac98,\qquad
  \beta=-\tfrac{15}{4},\qquad
  \beta-\alpha=-\sch
\]
exactly over $\mathbb{Q}$.
On every measured source window the same rewrite holds at
machine precision ($|\beta-\alpha-(-\sch)|\le 2\cdot 10^{-17}$)
for $\ell=v_{\mathrm{top}}$ and for $\ell=1$.
\end{lemma}

\begin{remark}
Lemma~\ref{lem:ab} is determinant / congruence bookkeeping
in the r401 Sylvester chart plus the r406 Woodbury form.
It does \emph{not} produce a new sign criterion:
$\beta>\alpha$ if and only if $\sch<0$, independently of
the bulk vector $\ell$ used to form $H$.
The closed formula
$\kappa=-(1-\eta)\det K_2/\det K_W$ of the $3\times 3$ toy
does not match $\kappa_{\mathrm{emp}}$ on the source two-column
split (that $U$ is not the $U$ of $\Phi$).
$\kappa=1$ is false: $\Delta=O(10^{-3})$ while
$|\sch|=O(10^{-2})$.
\end{remark}

\section{What the Euler tail sees on $v_{\mathrm{top}}$}

\begin{prop}[Euler re-read]\label{prop:euler}
Because $E_p+T_p=1$ pointwise,
$\sum_p\nu_p|\langle E_p+T_p,v\rangle|^2=|\langle 1,v\rangle|^2$
for every $v$.
At $k_z=9$, this equals $0.01192$ on $v_{\mathrm{top}}$ and
is \emph{not} $\alpha_T=-0.07302$.
The geometric Euler span misses $v_{\mathrm{top}}$
(QR-relres $0.74$).
A positive source representation of $\beta-\alpha$ by Euler
terms plus disk Parseval on $v_{\mathrm{top}}$ is therefore
not available.
The constant-vector control reproduces the r405 pins
($\Delta=2.110\cdot 10^{-3}$, $c_2=0.99789$,
$\kappa=0.03151$).
\end{prop}

\begin{prop}[not the defect]\label{prop:notdef}
$v_{\mathrm{top}}$ is not $\Tfrak_0$'s top singular vector
and not $A_0$'s negative eigenvector.
At $k_z=9$, $\cos=0.6766$ and $\alpha_T<0$ (contractive);
the true top singular value has Rayleigh $+0.16670$.
At $k_z=18$ (a genuine one-defect window,
$\|\Tfrak_0\|=1.25$) the cosine drops to $0.058$.
$H=A_0+v_{\mathrm{top}}v_{\mathrm{top}}^T$ stays indefinite
on $24/42$ MAIN windows (PD on only $18/42$).
The barycentric $Y$-top mode and the min-norm interpolant's
excess direction are different orthogonalities
(discrete OP on the holes versus dual OP chain of the full
measure, evaluated on $Y$).
\end{prop}

\section{Census and kills}

\begin{prop}[sign census]\label{prop:census}
On the named census, with the chart-scaled
$(\alpha,\beta)$ of \S\ref{sec:def}:
MAIN $42/42$ have $\beta>\alpha$ (all $\sch<0$);
living $\chi_3$ $37/37$ and living $\chi_4$ $41/41$
(together $78$) have $\beta>\alpha$;
dead $\chi_3$ $5/5$ and dead $\chi_4$-$20$ (together $6$)
have $\beta<\alpha$.
Death remains $\sch>0$ (r401), not a top-mode under-coverage
beyond that tautology.
Cosines $\cos(v_{\mathrm{top}},v_{-})$ range in
$[0.006,0.816]$.
\end{prop}

Kills, all of which fire.
(1)~The constant reproduces r405 (wrong vector, as designed).
(2)~Dropping $\sqrt{u^\vee}$ ($v\sim 1/P_Y'$) drops the
$k_z=9$ singular-vector cosine $0.68\to 0.38$.
(3)~Replacing $P_Y'$ by $y_i-\mathrm{mean}(Y)$ drops it to
$0.17$.
(4)~Weight-permute has $n_{-}(I-\Tfrak^*\Tfrak)=20$; scramble
has $21$.  The bulk is already dead; an edge balance of one
mode is then irrelevant.
(5)~Omitting the Sherman--Morrison border leaves the mixed
$2\times 2$ $K_2$ with $\indm=1$ on \emph{both} living $k_z=9$
and dead $\chi_3$-$15$; without the border there is no
living/dead scalar.

\section{What this is not}

This round does not prove HTM, QD, or P1.
It does not claim $R^\dagger\succ\tfrac12 I$.
It does not move the mincut
(\texttt{selected\_augDualResolvent\_gt\_half},
base $4$ / refined $5$).
The edge theorem asked for in DCCLXXIX, as a
\emph{top-mode-specific} positivity $\beta>\alpha$ with a
positive Euler source, is reduced to the r401 chart plus
the r405/r406 Woodbury rewrite.
Whether HTM itself holds is the parallel round 413.
No RH claim.

\end{document}
